\subsection{A masked sumcheck with hidden offsets}\label{sec:full:sumcheck}

Let $f:I_m\to\F$ be the private witness,
$z\in\F^m$ the public evaluation point,
and $v=\widetilde f(z)$ the authorized public
value. Put $\chi_0(a)=\operatorname{eq}_m(a,z)$
and
\[
 \sigma_0:=\sum_{a\in I_m}f(a)\chi_0(a)=v.
\]
For each round $1\le j\le m$ use the
unmasked degree-two polynomial $s_j(T)$
from~\eqref{eq:security:kbfold-round}.
Its challenge $r_{j-1}\in\F$ is sampled
only after the round-message commitment
has been bound.

\begin{definition}[True sumcheck states]\label{def:full:true-sum}
Starting with $(f_0,\chi_0)=(f,\chi_0)$,
put
\[
 f_j:=\Res_{r_{j-1}}f_{j-1},\qquad
 \chi_j:=\Res_{r_{j-1}}\chi_{j-1},
 \qquad
 \sigma_j:=\sum_{a\in I_{m-j}}f_j(a)\chi_j(a).
\]
The prover computes $s_j(T)$ from
$f_{j-1},\chi_{j-1}$ by
Definition~\ref{def:security:kbfold-transcript}.
\end{definition}

\begin{lemma}[Exact sumcheck transition]\label{lem:full:true-transition}
For every honest round $1\le j\le m$,
\[
 \boxed{
 s_j(0)+s_j(1)=\sigma_{j-1},\qquad
 s_j(r_{j-1})=\sigma_j.}
\]
\end{lemma}
\begin{proof}
The first equality is
Lemma~\ref{lem:security:kbfold-roundsum}.
For the second, each summand of
$s_j(r_{j-1})$ is the product of the
two affine restrictions
$\Res_{r_{j-1}}(f_{j-1})(b)$
and $\Res_{r_{j-1}}(\chi_{j-1})(b)$.
By Definition~\ref{def:full:true-sum}
these are $f_j(b)$ and $\chi_j(b)$.
Summing their products over
$b\in I_{m-j}$ gives $\sigma_j$.
\end{proof}

\begin{definition}[Hidden-offset Sumcheck]\label{def:full:masked-round}
Initialize $b_0:=0$. At the beginning
of round $j$, sample a fresh independent
$c_j\leftarrow\F$ and define
\[
 \begin{aligned}
 H_j(T)&:=\frac{b_{j-1}}{2}
              +c_j(2T-1),\\
 \mu_j(T)&:=s_j(T)+H_j(T),\\
 b_j&:=H_j(r_{j-1})
       =\frac{b_{j-1}}{2}
               +c_j(2r_{j-1}-1).
 \end{aligned}
\]
The prover sends \emph{only} a fresh,
perfectly hiding commitment
$\mathbf C_{\mu_j}$ to the three
coefficients of the degree-$\le2$
polynomial $\mu_j$. The coefficients,
$c_j$, $b_j$, and $s_j$ remain private.
The public challenge $r_{j-1}$ is
chosen after $\mathbf C_{\mu_j}$
is fixed.
\end{definition}

\begin{lemma}[Endpoint carry for the round mask]\label{lem:full:mask-endpoints}
For every $j$,
\[
 \boxed{H_j(0)+H_j(1)=b_{j-1}.}
\]
\end{lemma}
\begin{proof}
From the defining formula,
$H_j(0)=b_{j-1}/2-c_j$ and
$H_j(1)=b_{j-1}/2+c_j$.
Their sum is $b_{j-1}$.
The division by two is valid
because $\F$ has odd characteristic.
\end{proof}

\begin{theorem}[Exact masked sumcheck invariant]\label{thm:full:masked-sumcheck}
For every $1\le j\le m$, the
honest hidden round polynomial satisfies
\[
 \boxed{
 \begin{aligned}
 \mu_j(0)+\mu_j(1)&=\sigma_{j-1}+b_{j-1},\\
 \mu_j(r_{j-1})&=\sigma_j+b_j,
 \end{aligned}}
\]
and $\deg\mu_j\le2$.
Thus the publicly authorized
initial claim is
$\sigma_0+b_0=v$ and the
private running claim after
round $j$ is $\sigma_j+b_j$.
\end{theorem}
\begin{proof}
Lemma~\ref{lem:full:true-transition}
gives the two unmasked sumcheck
equalities. Lemma~\ref{lem:full:mask-endpoints}
gives the mask's endpoint sum,
and Definition~\ref{def:full:masked-round}
gives its evaluation at the
fresh challenge. Adding the
respective equalities proves
both displayed identities.
The degree bound holds because
$s_j$ has degree at most two
by Lemma~\ref{lem:security:kbfold-rounddegree}
and $H_j$ has degree at most one.
The starting relation follows
from $\sigma_0=v$ and $b_0=0$.
\end{proof}

\begin{lemma}[Public evaluation weight at the last round]\label{lem:full:final-chi}
At the end of all $m$ rounds,
the zero-variable weight is
the public scalar
\[
 \boxed{
 \chi_m=\prod_{i=0}^{m-1}
 \bigl((1-r_i)(1-z_i)+r_iz_i\bigr).
 }
\]
Moreover
$\sigma_m=\widetilde f(r)\chi_m$.
\end{lemma}
\begin{proof}
The function $\chi_0$ is the
Boolean equality table
$b\mapsto\operatorname{eq}_m(b,z)$.
Successive affine restrictions
at $r_0,\ldots,r_{m-1}$
produce the multilinear extension
$\widetilde{\chi_0}(r)$ by
Lemma~\ref{lem:kbff:prefix-mle}.
Since the equality table is the
product of coordinate-wise
Boolean weights, its
multilinear extension is the
displayed product.
The final table $f_m$ has
one entry, equal to
$\widetilde f(r)$ by
Corollary~\ref{cor:kbff:full-prefix}.
Its product with the single
entry $\chi_m$ is
$\sigma_m$ by definition.
\end{proof}

\begin{corollary}[Masked sumcheck to linear fold closure]\label{cor:full:masked-final-link}
Let $Q_m$ be the honestly masked
fold polynomial of
Theorem~\ref{thm:kbff:linear-closure},
and put $a:=[X^0]Q_m$. Then
\[
 \boxed{
 \mu_m(r_{m-1})=\chi_m a+b_m.
 }
\]
\end{corollary}
\begin{proof}
By Theorem~\ref{thm:full:masked-sumcheck},
$\mu_m(r_{m-1})=\sigma_m+b_m$.
Lemma~\ref{lem:full:final-chi}
makes $\sigma_m=\widetilde f(r)\chi_m$.
The constant-term closure identity
Corollary~\ref{cor:kbff:masked-constant}
gives $a=\widetilde f(r)$.
Substitution proves the equality.
\end{proof}

\begin{lemma}[Publicly revealing the masked round coefficients is unsafe]\label{lem:full:quadratic-leak}
The mask $H_j$ does not alter
the quadratic coefficient of
$s_j$. In particular,
\[
 [T^2]\mu_j=[T^2]s_j.
\]
Therefore the prescribed protocol
must not publish the entire
$\mu_j$ coefficient vector
and infer privacy merely from
the uniformity of $c_j$.
\end{lemma}
\begin{proof}
By Definition~\ref{def:full:masked-round},
$H_j$ has degree at most one.
Consequently its $T^2$
coefficient is zero.
Taking the coefficient of $T^2$
in $\mu_j=s_j+H_j$ proves
the displayed identity.
A public quadratic coefficient
would hence reveal the unmasked
quadratic coefficient exactly.
\end{proof}

\begin{remark}[Relation-proof requirement]
The verifier must certify the
endpoint sum, the masked challenge
transition, the updated private
offset $b_j$, and the final
equation of
Corollary~\ref{cor:full:masked-final-link}
\emph{without revealing} their
secret values. A generic zero-knowledge
proof of the explicit finite-field
relations will perform these checks;
ordinary raw sumcheck message
verification cannot.
\end{remark}
